\section*{Chapter \ref{ch:b1:stress-strain-rheology} --- Stress, Strain and Rheology}

\begin{solution}{exo:b1:stress-strain-rheology:1}
(a) $\sigma_v = 2700\times9.81\times10\,000 = \qty{2.65e8}{Pa} = \qty{264.9}{MPa}$.
(b) $P = 1000\times9.81\times10\,000 = \qty{98.1}{MPa}$. (c) $\sigma_v' =
264.9 - 98.1 = \qty{166.8}{MPa}$, and $\lambda = 98.1/264.9 = 0.370$, which is
$1000/2700$.
\end{solution}

\begin{solution}{exo:b1:stress-strain-rheology:2}
(a) $e = (13-10)/10 = 0.30$ and $S = 1.30$. (b) $e = (7-10)/10 = -0.30$, a
shortening of \qty{30}{\percent}, and $S = 0.70$. (c) The angle has closed by
$\psi = \ang{30}$, and $\gamma = \tan\ang{30} = 0.577$.
\end{solution}

\begin{solution}{exo:b1:stress-strain-rheology:3}
(a) Centre \qty{70}{MPa}, radius \qty{30}{MPa}. (b) The largest shear stress
is the radius, \qty{30}{MPa}, on the planes at \ang{45} to $\sigma_1$ and
$\sigma_3$, which carry the \omterm{def:b1:stress-strain-rheology:stress}{normal stress} of the centre, \qty{70}{MPa}. (c)
$2\theta = \ang{60}$: $\sigma_n = 70 + 30\cos\ang{60} = \qty{85}{MPa}$ and
$\tau = 30\sin\ang{60} = \qty{26.0}{MPa}$.
\end{solution}

\begin{solution}{exo:b1:stress-strain-rheology:4}
(a) $\ang{45} - \ang{30}/2 = \ang{30}$ to $\sigma_1$, on either side of it,
so the two planes make \ang{60} with each other. (b) With $c' = 20 + r$, the
condition of \cref{thm:b1:stress-strain-rheology:failure-angle} reads
$r = C\cos\phi + (20 + r)\sin\phi$, so
\[
  r = \frac{C\cos\phi + 20\sin\phi}{1 - \sin\phi} = \frac{8.66 + 10}{0.5}
    = \qty{37.3}{MPa},
\]
and $\sigma_1' = 20 + 2r = \qty{94.6}{MPa}$.
\end{solution}

\begin{solution}{exo:b1:stress-strain-rheology:5}
(a) $c = 40$ and $R = \sqrt{20^2 + 15^2} = \qty{25}{MPa}$: $\sigma_1 =
\qty{65}{MPa}$ and $\sigma_3 = \qty{15}{MPa}$; $\tan2\theta_p = 30/40 = 0.75$,
$2\theta_p = \ang{36.87}$, $\theta_p = \ang{18.43}$. (b) At this angle
$\cos^2\theta = 0.9$, $\sin^2\theta = 0.1$ and $\sin\theta\cos\theta = 0.3$:
$\sigma_n = 60\times0.9 + 2\times15\times0.3 + 20\times0.1 = 54 + 9 + 2 =
\qty{65}{MPa}$, and $\tau = (20-60)\times0.3 + 15\times(0.9-0.1) = -12 + 12 =
0$. (c) $R = \qty{25}{MPa}$, on the planes at \ang{45} to the principal
directions.
\end{solution}

\begin{solution}{exo:b1:stress-strain-rheology:6}
(a) $S_{1,2} = \sqrt{1 + 1} \pm 1$: $S_1 = 2.414$ and $S_2 = 0.414$;
$\tan2\theta' = 2/2 = 1$, $\theta' = \ang{22.5}$. (b) $S_1S_2 =
(\sqrt2 + 1)(\sqrt2 - 1) = 1$, which is $\det F$: no change of area. (c) The
line along $(1, 1)$ goes to $F(1, 1) = (1 + 2, 1) = (3, 1)$, at
$\arctan(1/3) = \ang{18.4}$ from the shear direction, with the \omterm{def:b1:stress-strain-rheology:strain}{stretch}
$\sqrt{10}/\sqrt2 = 2.24$. It has turned past the long axis of the ellipse,
at \ang{22.5}, towards the shear plane; its \omterm{def:b1:stress-strain-rheology:strain}{stretch} is a little less than
$S_1$, as it must be, since no line \omterm{def:b1:stress-strain-rheology:strain}{stretches} more than $S_1$.
\end{solution}

\begin{solution}{exo:b1:stress-strain-rheology:7}
(a) $S_1/S_2 = 4$ and $S_1S_2 = 1$ give $S_1 = 2$ and $S_2 = 0.5$; by
\cref{prop:b1:stress-strain-rheology:simple-shear}, $\gamma = S_1 - S_2 =
1.5$. (b) $\tan2\theta' = 2/1.5 = 1.333$, $2\theta' = \ang{53.1}$: the long
axes should lie at \ang{26.6} to the walls. (c) $1.5\times200 = \qty{300}{m}$.
\end{solution}

\begin{solution}{exo:b1:stress-strain-rheology:8}
$\sigma_v = 2700\times9.81\times3000 = \qty{79.46}{MPa}$. (a) $P = \qty{29.43}{MPa}$,
$\sigma_n' = 70 - 29.43 = \qty{40.57}{MPa}$, and Byerlee allows
$0.85\times40.57 = \qty{34.5}{MPa}$, more than the \qty{30}{MPa} applied: the
fault holds. (b) $P = 0.6\times79.46 = \qty{47.68}{MPa}$, $\sigma_n' =
\qty{22.32}{MPa}$, and the friction gives only \qty{19.0}{MPa}: the fault
slips. (c) $30 = 0.85\,(70 - 79.46\,\lambda)$ gives $\lambda = (70 -
35.29)/79.46 = 0.437$.
\end{solution}

\begin{solution}{exo:b1:stress-strain-rheology:9}
(a) The stress goes as $\dot\varepsilon^{1/4}$: $10^{1/4} = 1.78$. (b) The
ratio of the strengths is
\[
  \exp\Bigl[\frac{E}{nR}\Bigl(\frac{1}{673.15} - \frac{1}{773.15}\Bigr)\Bigr]
  = \exp\bigl(6705\times\num{1.921e-4}\bigr) = \exp(1.288) = 3.63 .
\]
(c) A tenfold change of rate changes the creep strength by a factor 1.78,
which the temperature does in only $100\times\ln1.78/\ln3.63 \approx
\qty{45}{\degreeCelsius}$; on a \omterm{def:b1:heat-budget:geotherm}{geotherm} of about \qty{20}{\degreeCelsius}
per kilometre, that moves the transition by about \qty{2}{km}. A change of
\omterm{def:b1:heat-budget:geotherm}{geotherm} moves the isotherms, and with them the transition, by far more.
\end{solution}

\begin{solution}{exo:b1:stress-strain-rheology:10}
(a) With $C = 0$, a fault slips when $r = \mu c'/\sqrt{1+\mu^2}$
(\cref{thm:b1:stress-strain-rheology:failure-angle}). Here $c' = S =
\rho g z(1-\lambda)$, so $\sigma_1 - \sigma_3 = 2r = 2\mu S/\sqrt{1+\mu^2}$.
(b) $S$ grows by $2700\times9.81\times1000\times(1 - 1000/2700) =
\qty{16.68}{MPa}$ per kilometre, and $2\times0.85/1.3124 = 1.295$: the slope
is \qty{21.6}{MPa} per kilometre, between \num{13.1} (normal faults) and
\qty{61.3}{MPa} per kilometre (thrusts). (c) $2\mu/(k+\mu) = 0.786$,
$2\mu/k = 1.295$ and $2\mu/(k-\mu) = 3.68$, with $k = 1.3124$: the
\omterm{def:b1:stress-strain-rheology:principal}{differential stress} at slip is 0.79, 1.30 and 3.68 times the vertical
\omterm{def:b1:stress-strain-rheology:lithostatic}{effective stress}.
\end{solution}

\begin{solution}{exo:b1:stress-strain-rheology:11}
(a) $t = 60\times\num{e6}\times\num{3.156e7} = \qty{1.89e15}{s}$ and
$2\sqrt{\kappa t} = 2\sqrt{\num{1.89e9}} = \qty{87.0}{km}$. (b)
$\operatorname{erf}x = 600/1350 = 0.444$ for $x = 0.417$: $z = 0.417\times87.0
= \qty{36.3}{km}$; $\operatorname{erf}x = 800/1350 = 0.593$ for $x = 0.586$:
$z = \qty{51.0}{km}$. (c) The earthquakes should reach \qtyrange{36}{51}{km};
the envelope puts the base of the brittle layer at \qty{38.2}{km}, within
that range and near its top, at a temperature, \qty{641}{\degreeCelsius}, in
the range observed.
\end{solution}

\begin{solution}{exo:b1:stress-strain-rheology:12}
(a) $r = 45$ and $c' = \qty{65}{MPa}$; the distance from the centre to the line
$\tau = 0.85\sigma_n'$ is $0.85\times65/1.3124 = \qty{42.1}{MPa}$, less than
the radius: the circle crosses the line, and the planes on the arc above it
slip. (b) With $2\theta$ the angle from the $\sigma_1$ end, slip needs
$45\sin2\theta \ge 0.85\,(65 + 45\cos2\theta)$, that is
$45\sin2\theta - 38.25\cos2\theta \ge 55.25$, or
$59.06\sin(2\theta - \ang{40.36}) \ge 55.25$, where $\ang{40.36} = \phi$.
So $\sin(2\theta - \phi) \ge 0.9355$, $2\theta - \phi$ between \ang{69.31}
and \ang{110.69}, $2\theta$ between \ang{109.67} and \ang{151.06}: the normals
lie at \ang{54.8} to \ang{75.5} from $\sigma_1$, and the slipping fractures
at \ang{14.5} to \ang{35.2} from $\sigma_1$. The middle of the range,
\ang{24.8}, is $\ang{45} - \phi/2$, the plane of
\cref{thm:b1:stress-strain-rheology:failure-angle}. (c) Only that plane slips
when the circle just touches the line: $r = \mu(20 + r)/k$, $r = 20\mu/(k -
\mu) = 17/0.4624 = \qty{36.8}{MPa}$, and $\sigma_1' = 20 + 2r =
\qty{93.5}{MPa}$.
\end{solution}

\begin{solution}{pb:b1:stress-strain-rheology:1}
\textbf{1.} $\sigma_v = 2700\times9.81\times5000 = \qty{132.44}{MPa}$.

\textbf{2.} $P = 1000\times9.81\times5000 = \qty{49.05}{MPa}$, and
$\lambda = 49.05/132.44 = 0.370$.

\textbf{3.} $\sigma_1 = \sigma_v = \qty{132.44}{MPa}$, vertical, and
$\sigma_3 = \qty{75.00}{MPa}$, horizontal east--west: the \omterm{def:b1:layered-earth:layers}{crust} is being
stretched east--west under its own weight, the state that drives normal
faults, such as this one.

\textbf{4.} $\sigma_1 - \sigma_3 = \qty{57.44}{MPa}$; centre
$(132.44 + 75.00)/2 = \qty{103.72}{MPa}$; radius \qty{28.72}{MPa}.

\textbf{5.} $\sigma_1' = 132.44 - 49.05 = \qty{83.39}{MPa}$, $\sigma_3' =
75.00 - 49.05 = \qty{25.95}{MPa}$, centre $103.72 - 49.05 = \qty{54.67}{MPa}$.

\textbf{6.} $\sigma_{xx} = \qty{75.00}{MPa}$, $\sigma_{yy} = \qty{132.44}{MPa}$
and $\sigma_{xy} = 0$: the axes are the principal directions, one of them
vertical, so the planes perpendicular to them carry no shear stress.

\textbf{7.} Down its \omterm{def:b1:geological-maps:attitude}{dip} the fault runs along the unit vector
$(\cos\ang{60}, -\sin\ang{60})$, east and down; its upward normal is
$\vec n = (\sin\ang{60}, \cos\ang{60}) = (0.866, 0.500)$, at \ang{30} from
$x$. Its angle with $\sigma_1$, along $y$, has the cosine
0.500: $\theta = \ang{60}$.

\textbf{8.} With \ang{30} from $x$: $\sigma_n = 75.00\times0.75 + 132.44\times0.25
= 56.25 + 33.11 = \qty{89.36}{MPa}$, and $\tau = (132.44 - 75.00)\times
\sin\ang{30}\cos\ang{30} = 57.44\times0.4330 = \qty{24.87}{MPa}$.

\textbf{9.} At $2\theta = \ang{120}$ from the $\sigma_1$ end:
$\sigma_n = 103.72 + 28.72\cos\ang{120} = 103.72 - 14.36 = \qty{89.36}{MPa}$
and $\tau = 28.72\sin\ang{120} = \qty{24.87}{MPa}$.

\textbf{10.} $\sigma_n' = 89.36 - 49.05 = \qty{40.31}{MPa}$.

\textbf{11.} $24.87/40.31 = 0.617$, below 0.85: the fault holds.

\textbf{12.} $0.85\times40.31 = \qty{34.26}{MPa}$; the shear stress falls
short of it by \qty{9.39}{MPa}.

\textbf{13.} The circle has its centre at \qty{54.67}{MPa} and the radius
\qty{28.72}{MPa}, from \num{25.95} to \qty{83.39}{MPa}; Byerlee's line is
$\tau = 0.85\,\sigma_n'$ (all normal stresses here are below
\qty{200}{MPa}); the fault is the point $(40.31, 24.87)$, at \ang{120} from
the $\sigma_1$ end.

\textbf{14.} $0.85\times54.67/\sqrt{1 + 0.85^2} = 46.47/1.3124 =
\qty{35.41}{MPa}$, more than the radius, \qty{28.72}{MPa}: the whole circle
lies below the line, and no plane through the rock can slip.

\textbf{15.} $\phi = \arctan0.85 = \ang{40.36}$: the plane at
$\ang{45} - \phi/2 = \ang{24.8}$ from the vertical $\sigma_1$, dipping
\ang{65.2}.

\textbf{16.} $S = \sigma_v' = \qty{83.39}{MPa}$ and
$\sigma_1 - \sigma_3 = 2\times0.85\times83.39/(1.3124 + 0.85) =
\qty{65.55}{MPa}$; the slope gives $13.11\times5 = \qty{65.55}{MPa}$.

\textbf{17.} $\sigma_3 = 132.44 - 65.55 = \qty{66.88}{MPa}$.

\textbf{18.} The fault slips when
$r\sin\ang{120} = 0.85\,(c' + r\cos\ang{120})$ with
$c' = (83.39 + \sigma_3')/2$ and $r = (83.39 - \sigma_3')/2$; solving gives
$\sigma_3' = \qty{17.18}{MPa}$, $\sigma_3 = \qty{66.23}{MPa}$. It is lower
than for the best plane: the fault does not lie at the point where the circle
first touches the line, so the circle must grow further before it reaches
the fault's own point.

\textbf{19.} Every effective \omterm{def:b1:stress-strain-rheology:stress}{normal stress}, $\sigma_1'$ and $\sigma_3'$
included, falls by $\Delta P$, and their difference does not change; the
shear stresses do not depend on the pore pressure at all.

\textbf{20.} The best plane slips when the distance equals the radius:
$0.85\,(54.67 - \Delta P)/1.3124 = 28.72$, so
$\Delta P = 54.67 - 28.72\times1.3124/0.85 = 54.67 - 44.34 = \qty{10.33}{MPa}$.

\textbf{21.} The column of water in the well weighs the same as before, so the
pump must add the whole \qty{10.33}{MPa} at the wellhead, about 103 times
the atmospheric pressure: the weight of a column of water
$\num{10.33e6}/(1000\times9.81) = \qty{1053}{m}$ high.

\textbf{22.} With $\mu = 0.6$ the fault could bear only
$0.6\times40.31 = \qty{24.19}{MPa}$, less than the \qty{24.87}{MPa} it
carries: it would be slipping already. Since it is quiet, its friction is at
least $0.617$, or it has some \omterm{def:b1:stress-strain-rheology:coulomb}{cohesion}.

\textbf{23.} The best-oriented plane slips when the pore pressure reaches
$49.05 + 10.33 = \qty{59.38}{MPa}$, that is $\lambda = 59.38/132.44 = 0.448$,
well short of the lithostatic value 1.

\textbf{24.} $\Delta P = \sigma_n' - \tau/0.85 = 40.309 - 24.870/0.85 = 40.309 -
29.259 = \qty{11.049}{MPa}$, more than the \qty{10.33}{MPa} that would move a
plane of the best orientation. To one decimal, the rise of pore pressure that
brings the fault to slip is
\textbf{\qty[text-series-to-math=true]{11.0}{MPa}}.
\end{solution}
